% ===========================================================================
%  9 — Contratos de integracao
% ===========================================================================
\chapter{Contratos de integração}
\label{cap:contratos}

Há quatro fronteiras nesta plataforma cuja quebra é caríssima, porque atinge
código que já está publicado e não pode ser atualizado em conjunto: o protocolo
\cd{postMessage} entre hub e experiência, o formato do evento de telemetria, o
manifesto de experiências e a interface do \emph{provider} de IA. Este capítulo
especifica os quatro e a política de versionamento de cada um --- e, desde
2026-09-30, um quinto, que tem quatro pontas em vez de duas: o do agente de IA do
professor (Seção~\ref{sec:contrato-agente}). Desde 2026-10-02 há um sexto, o
único em que quem chama a api é o servidor de um \emph{vendor}: o do avatar da
Ajuda (Seção~\ref{sec:contrato-ajuda}).

\begin{regra}
\textbf{Mudança em contrato exige cuidado extremo e atualização deste
documento.} Um \emph{bundle} publicado há três meses no \emph{bucket} de um
contrato continua sendo servido a estudantes: ele não é reconstruído quando a API
muda. Toda alteração de contrato é, por definição, uma alteração
\emph{retrocompatível} ou um \emph{bump} de versão maior com suporte à anterior.
\end{regra}

\section{Protocolo \cd{postMessage} hub $\leftrightarrow$ experiência}
\label{sec:bridge}

\subsection{Especificação das mensagens (protocolo v1)}

\begin{tabularx}{\linewidth}{@{}L{18mm}L{34mm}X@{}}
\toprule
\headrow \thd{Sentido} & \thd{Mensagem} & \thd{Carga e semântica}\\
\midrule
exp \seta{} hub & \cd{exp:ready} &
  \cd{\{ protocolVersion \}}. Anúncio de que o \emph{bundle} carregou. É
  \textbf{reenviado} a cada \SI{400}{\milli\second} até o hub responder.\\
\zebra hub \seta{} exp & \cd{hub:init} &
  \cd{\{ protocolVersion, token, estudante\{id,nome,serie,turmaId,escolaId\},
  apiUrl, config \}}. Resposta ao \cd{exp:ready}; só a primeira conta.
  \cd{config} é o saco aditivo do protocolo: hoje carrega \cd{expId},
  \cd{clarityProjectId} (\adrr{ADR-60}), \cd{obs.traceRatio} (a fração de
  sessões com OpenTelemetry Web, Capítulo~\ref{cap:observabilidade}) e
  \cd{fuso} (o fuso da plataforma, \adrr{ADR-99}) --- campo novo aqui não muda
  \cd{PROTOCOL\_VERSION}, porque a experiência que não o conhece o ignora.
  Os três são valores de \emph{runtime} do hub, e nenhum poderia ser assado no
  \emph{bundle}: ele é construído uma vez e não conhece contrato.\\
exp \seta{} hub & \cd{exp:concluir} &
  \cd{\{ payload: \{ score?, dados? \} \}}. Conclusão da missão. O hub registra
  progresso.\\
\zebra exp \seta{} hub & \cd{exp:token-expirando} &
  Sem carga. Pedido de renovação; o hub responde \cd{hub:token}.\\
hub \seta{} exp & \cd{hub:token} &
  \cd{\{ token \}}. Novo \emph{token}, sem recarregar o \cd{iframe}.\\
\zebra exp \seta{} hub & \cd{exp:estado} &
  \cd{\{ estado: \{\ldots\} \}}. Estado de alto nível para o \emph{chrome}
  opcional do hub. \textbf{Aditivo} ao v1; o hub \emph{não} interpreta domínio,
  só repassa a quem plugou um \emph{callback}.\\
exp \seta{} hub & \cd{exp:erro} & \cd{\{ mensagem, hash \}}: a experiência
  quebrou de um jeito que ela não trata (\cd{window.onerror},
  \cd{unhandledrejection}). Só para o hub TROCAR A TELA (\enquote{A missão
  travou}, recarregar); o registro do erro vai da SDK direto à API
  (\rota{POST /observabilidade/erros-cliente}, \adrr{ADR-58}). Aditivo, v2
  continua; hub antigo descarta com o aviso de \emph{schema}.\\
hub \seta{} exp & \cd{hub:pause} / \cd{hub:resume} &
  Sem carga. Emitidos por visibilidade e foco da página do hub.\\
\bottomrule
\end{tabularx}

\begin{nota}[title={NOTA --- o que o hub faz com o \cd{exp:estado}}]
O \cd{ExpHost} guarda o último estado publicado e reconhece dois campos, ambos
da exp08ef1b: \cd{acao: "sair"} (a sessão é do HUB, em \emph{cookie} \cd{httpOnly}
fora do alcance do \emph{iframe}, então a experiência pede e o hub executa) e
\cd{fase}, que acende a moldura dos posts perdidos e --- com
\cd{passo: "relatorio"} --- solta o aparelho para o relatório ocupar a área
inteira (ADR-41, \ref{adr:relatorio-expandido}).

Duas consequências de o hub SUBSTITUIR o estado inteiro a cada mensagem:
publicar só o \cd{passo} apagaria a \cd{fase}, então quem publica leva as duas;
e é a própria troca de fase seguinte que desfaz o estado anterior, sem
mensagem de ``desligar''. As outras experiências publicam estado
(\cd{\{ tela \}} na exp04ef1b, \cd{\{ passo, tentativa \}} na exp02em1b-v0) e o hub
ignora --- o canal é genérico de propósito. O segundo \emph{chrome} é a
\textbf{instrução ao lado do aparelho} (exp02em1b 2.1.0, \adrr{ADR-191}):
\cd{\{ instrucao: \{ antes, destaque, depois \} \}}, validado com Zod, desenhado
pelo hub à direita do \cd{DispositivoFrame} (acima dele em tela estreita), com
o \cd{destaque} em negrito; sem \cd{instrucao}, nada aparece. O hub não sabe de
que missão é o texto --- é o TC-10, com o tema do professor, que a api manda
e a experiência publica nas telas do feed.
\end{nota}

\subsection{O \emph{handshake}, passo a passo}

\begin{figure}[H]
\centering
\fitwidth{%
\begin{tikzpicture}[x=1cm, y=1cm]

  % Raias: experiência, ExpFrame, rota do hub e api. Rótulo longo quebra em
  % duas linhas para caber no vão entre as raias (nenhum cruza uma linha de vida).
  \raia{1.4}{-15.9}{Experiência}{\texttt{@edu/sdk} no \texttt{iframe}}
  \raia{6.2}{-15.9}{\texttt{ExpFrame}}{React, no hub}
  \raia{10.4}{-15.9}{Rota do hub}{Next.js, servidor}
  \raiaext{15.0}{-15.9}{API}{NestJS}

  % 1. A rota do hub monta o embed (servidor)
  \msgs{10.4}{15.0}{-1.1}{\texttt{GET /experiencias/:expId}}
  \msgr{15.0}{10.4}{-2.1}{manifesto \texttt{\{ versaoAtiva,}\\\texttt{viewport, permissoes \}}}
  \msgself{10.4}{-2.6}{assina o \emph{token}\\RS256, 15 min}
  \msgs{10.4}{6.2}{-4.0}{\emph{props}: \texttt{src}, \emph{token}, \texttt{apiUrl}}

  % 2. Handshake no navegador
  \msgs{6.2}{1.4}{-5.0}{renderiza o \texttt{iframe}\\com a URL versionada}
  \msgself{1.4}{-5.5}{\texttt{connectToHub}\\\texttt{(\{ hubOrigin \})}}
  \draw[seqframe] (0.8,-6.3) rectangle (6.8,-8.0);
  \node[frameop] at (0.8,-6.3) {loop};
  \node[guarda] at (1.85,-6.62) {[a cada 400 ms]};
  \msgs{1.4}{6.2}{-7.75}{\texttt{exp:ready}\\\texttt{\{ protocolVersion: 2 \}}}
  \msgself{6.2}{-8.35}{valida \texttt{ev.origin}\\e a versão MAIOR}
  \msgr{6.2}{1.4}{-9.9}{\texttt{hub:init \{ token, estudante,}\\\texttt{apiUrl, config \}}}
  \msgself{1.4}{-10.4}{valida com Zod;\\resolve a \emph{Promise}}

  % 3. Uso e renovação do token
  \draw[seqmsg] (1.4,-12.3) -- (15.0,-12.3)
    node[seqlbl, pos=0.24, above=1pt, fill=white, inner sep=1.5pt]{chamadas de domínio,\\\texttt{/telemetry/batch},\\\texttt{/ai/gerar}};
  \msgs{1.4}{6.2}{-13.3}{\texttt{exp:token-expirando}}
  \msgs{6.2}{10.4}{-14.1}{\texttt{GET /api/exp-token}}
  \msgr{10.4}{6.2}{-14.8}{\texttt{\{ token \}}}
  \msgr{6.2}{1.4}{-15.5}{\texttt{hub:token \{ token \}}}

  % Notas: abaixo das raias, lado a lado, para o desenho caber em 16 cm
  \node[actnote, text width=72mm, anchor=north west] at (0.1,-16.3)
    {\textbf{Por que o reenvio do \texttt{exp:ready}:} com os \emph{assets} em
     \emph{cache} (um F5), a experiência fica pronta ANTES de o \texttt{ExpFrame}
     hidratar o \emph{listener}, e um disparo único se perderia --- o
     \emph{handshake} expiraria sem culpa de ninguém. O hub responde a cada
     \texttt{exp:ready} com um \texttt{hub:init}; só o primeiro conta, e
     \texttt{hub:init} repetido é tratado como evento normal, não como ruído.};

  \node[actnote, text width=72mm, anchor=north west] at (8.2,-16.3)
    {\textbf{Por que a renovação passa pelo hub:} a experiência não tem sessão
     --- ela tem um \emph{token} de 15 minutos. Quem tem sessão é o hub, e é ele
     que pode assinar outro. O \texttt{iframe} nunca recarrega, então o estado da
     missão sobrevive à renovação.};

\end{tikzpicture}}
\caption[Sequência do handshake hub-experiência]{Sequência do \emph{handshake} e
  da renovação de \emph{token}. O \cd{targetOrigin} é explícito nos dois
  sentidos, e as duas pontas descartam mensagem de \emph{origin} divergente
  antes de olhar o conteúdo.}
\label{fig:seq-handshake}
\end{figure}

\subsection{Regras invioláveis do protocolo}

\begin{regra}
\begin{enumerate}
  \item \textbf{\cd{targetOrigin} explícito, nunca \cd{"*"}.} O hub usa a
    \emph{origin} derivada da \emph{mesma} URL que colocou no \cd{iframe}; a SDK
    usa o \cd{hubOrigin} que recebeu por \emph{query string}. Com \cd{"*"}, o
    \emph{token} iria para qualquer documento que conseguisse virar o pai do
    \emph{frame}.
  \item \textbf{Aquela \emph{query string} é entrada NÃO CONFIÁVEL, e passa por
    \cd{normalizarHubOrigin()}.} O \emph{bundle} é público: quem embeda a
    experiência num \cd{iframe} próprio escreve o parâmetro, e com ele escolheria
    de quem se aceita o \cd{hub:init} e para onde vai tudo o que sai depois. A
    função reduz o valor a uma origem canônica, recusa esquema fora de http(s),
    credencial e origem opaca, exige hub em https quando a experiência está em
    https, e confere contra \cd{location.ancestorOrigins[0]} --- a origem REAL de
    quem embedou, que VENCE o parâmetro (\adrr{ADR-46}).
  \item \textbf{Toda mensagem pós-\emph{init} é validada com Zod nos DOIS
    lados.} Mensagem inválida é descartada com aviso, nunca despachada crua aos
    \emph{listeners}.
  \item \textbf{Versão MAIOR verificada nos dois lados; incompatível falha
    alto.} O hub não envia \cd{hub:init} e registra erro no console; a
    experiência expira o \emph{handshake} e mostra a própria mensagem de erro.
  \item \textbf{O \emph{token} vive só em memória} na experiência. Nunca
    \cd{localStorage}, nunca \emph{cookie}.
  \item \textbf{Nunca quebrar compatibilidade sem \emph{bump} de versão maior e
    suporte à anterior.}
\end{enumerate}
\end{regra}

\begin{lstlisting}[language=tsunoo, caption={O núcleo da validação, em \arq{packages/sdk/src/bridge.ts}}, label={lst:bridge}]
function handler(ev: MessageEvent) {
  if (ev.origin !== hubOrigin) return;           // regra 1: valida origin

  if (!conectado) {
    const init = HubInitSchema.safeParse(ev.data);
    if (!init.success) return;                   // antes do init, tudo mais e ruido
    // Versao MAIOR precisa bater — hub e experiencia de geracoes diferentes
    // falham ALTO, nao silenciosamente.
    if (Math.trunc(init.data.protocolVersion) !== PROTOCOL_VERSION) {
      reject(new Error(`Protocolo incompativel: hub v${init.data.protocolVersion}, ` +
                       `experiencia v${PROTOCOL_VERSION}`));
      return;
    }
    conectado = true;
    resolve({ init: init.data, send, onMessage });
    return;
  }

  // hub:init repetido (resposta a um exp:ready reenviado) nao e ruido nem erro.
  if ((ev.data as { type?: unknown })?.type === "hub:init") return;

  const msg = HubParaExpSchema.safeParse(ev.data);
  if (!msg.success) {
    console.warn("[sdk] mensagem do hub descartada (schema):", ev.data?.type);
    return;                                      // regra 2: descarta, nao despacha
  }
  listeners.forEach((cb) => cb(msg.data));
}
\end{lstlisting}

\subsection{Política de evolução}

\begin{tabularx}{\linewidth}{@{}L{34mm}L{22mm}X@{}}
\toprule
\headrow \thd{Tipo de mudança} & \thd{Versão} & \thd{Regra}\\
\midrule
Nova mensagem opcional (como \cd{exp:estado}) & v1, aditivo &
  Permitido sem \emph{bump}: hub antigo ignora mensagem que não conhece, e
  experiência antiga nunca recebe a nova sem pedir.\\
\zebra Novo campo opcional em \cd{hub:init} & v1, aditivo &
  Permitido. O esquema Zod tem \cd{.default()} ou \cd{.optional()}, e a
  experiência antiga simplesmente não o lê.\\
Campo existente muda de tipo ou semântica & \textbf{v2} &
  Exige \emph{bump} de versão maior \emph{e} que o hub continue respondendo v1
  enquanto houver \emph{bundle} v1 publicado em algum contrato.\\
\zebra Remoção de mensagem ou de campo & \textbf{v2} &
  Mesma regra. A remoção só é efetiva quando o último \emph{bundle} v1 tiver sido
  republicado.\\
\bottomrule
\end{tabularx}

\begin{atencao}
\cd{hub:pause} e \cd{hub:resume} existem porque \textbf{esconder um
\cd{iframe} não dispara \cd{visibilitychange} nele}. Sem essas mensagens, o
cronômetro de \emph{dwell} da experiência correria com o estudante de costas --- e
\emph{dwell} é a medida central do motor de \emph{feed} da exp08ef1b. O hub escuta
\cd{visibilitychange} \emph{e} \cd{blur}/\cd{focus} da janela, porque trocar de
janela ou de aplicativo não muda \cd{visibilityState}.

\textbf{Uma exceção, e ela custou um bug:} clicar em qualquer lugar
\textbf{dentro} do \cd{iframe} move o foco para a experiência, e o navegador
dispara \cd{blur} no documento de cima. Isso \emph{não} é o estudante saindo --- e
como o \cd{focus} correspondente não volta enquanto ele continua ali, um
\cd{hub:pause} nesse caso é definitivo. Na exp04ef1b o vídeo parava no primeiro
toque em \enquote{Salvar} e não voltava mais. Quando o foco foi para a nossa
experiência, \cd{document.activeElement} é o próprio \cd{iframe}, e
\arq{ExpFrame.tsx} não emite nada.
\end{atencao}

\section{Manifesto de experiências}
\label{sec:manifesto}

O manifesto é o contrato entre o CI, a API e o hub. Ele existe em duas cópias
que devem concordar: o arquivo \arq{public/experiencia.json} dentro do
\emph{bundle}, e a linha da tabela \cd{experiencias} na API --- que o CI escreve
lendo o arquivo.

\begin{lstlisting}[language=jsonunoo, caption={\arq{apps/experiencias/app-70-vila-do-lume/public/experiencia.json}}, label={lst:manifesto}]
{
  "id": "app-70-vila-do-lume",
  "nome": "Vila do Lume — A Noite Longa",
  "versao": "1.1.3",
  "produto": "entre-redes",
  "series": [],
  "capabilities": ["telemetria"],
  "permissoes": ["pointer-lock"],
  "verbos": {
    "attempted":  "entrou na vila (partida aberta)",
    "interacted": "acendeu fogueira, gastou tocha, gravou na Pedra dos Nomes ou respondeu a uma mensagem",
    "progressed": "acendeu mais uma das 8 fogueiras",
    "completed":  "acendeu as 8 fogueiras e reacendeu o Fogo Sagrado"
  }
}
\end{lstlisting}

\begin{tabularx}{\linewidth}{@{}L{26mm}L{22mm}X@{}}
\toprule
\headrow \thd{Campo} & \thd{Tipo} & \thd{Semântica}\\
\midrule
\cd{id} & \emph{string} & Slug da experiência; é a chave primária da tabela e o
  primeiro segmento do caminho no \emph{bucket}.\\
\zebra \cd{nome} & \emph{string} & Nome de exibição.\\
\cd{versao} & \emph{semver} & Vem do \arq{package.json} da experiência. É o
  segundo segmento do caminho no \emph{bucket}, o que torna o \emph{bundle}
  imutável.\\
\zebra \cd{series} & \emph{string[]} & Séries em que a experiência se aplica.
  \textbf{Vazio significa TODAS as séries}, não nenhuma --- ver o alerta
  abaixo.\\
\cd{unidade} & \emph{1--4?} & A unidade do ano a que a missão pertence. Ausente
  na transversal e na ferramenta de professor: elas não moram numa unidade. É o
  que ORDENA a grade da sala (\adrr{ADR-69}).\\
\zebra \cd{produto} & \emph{slug} & A que PRODUTO a experiência pertence:
  \cd{entre-redes}, \cd{plano-de-aula}, \cd{plano-de-ensino} ou
  \cd{bem-estar}. É o que a põe no painel da Unooria e o que lhe dá o TETO do
  produto (\adrr{ADR-87}) --- e é por ele que a composição produto \seta{}
  experiências é DERIVADA do manifesto, em vez de ser uma lista mantida à mão
  que ninguém atualiza quando a experiência seguinte nascer. Cobrado por
  \arq{tests/manifestos-declaram-produto.spec.ts}, e \textbf{não} pelo Zod da
  rota: um 400 no \rota{PUT /experiencias/:expId} mataria a ativação DEPOIS de
  os \emph{bundles} já terem subido.\\
\cd{capabilities} & \emph{string[]} & \cd{telemetria}, \cd{progresso},
  \cd{ia}, \cd{avaliacao}, \cd{tempo-real}. Declara o que a experiência usa da
  plataforma.\\
\zebra \cd{viewport} & \cd{\{largura, altura\}?} & Presente \implica{} o hub
  embeda dentro do \cd{DispositivoFrame} (casca de aparelho com escala). A
  exp08ef1b e a exp06ef2b são 375$\times$812; a exp04ef1b \emph{não} tem \cd{viewport},
  porque foi desenhada para \emph{desktop}. Quem desenha o aparelho é o hub ---
  a experiência \textbf{não} empilha uma casca própria.\\
\cd{publico} & \cd{estudante|professor} & \cd{professor} \implica{} ferramenta do
  professor, e \rota{/exp/*} responde \textbf{404} para ela.\\
\zebra \cd{permissoes} & \emph{string[]} & Permissões que o \emph{iframe} daquela
  experiência precisa: \cd{pointer-lock} (app-70), \cd{forms} (exp02em1b, exp04ef1b, exp05ef1b, exp06ef1b, exp07ef1b, exp08ef1b e exp08ef4b) e
  \cd{downloads} (app-90, app-92), traduzidos pelo \cd{ExpFrame} nos tokens
  \cd{allow-pointer-lock}, \cd{allow-forms} e \cd{allow-downloads}. Vazio (o
  \emph{default}) = nenhuma. \textbf{O vocabulário é fechado em três lugares}
  --- o Zod da api, o CHECK de \cd{experiencias.permissoes} e o \cd{ExpFrame}
  ---, e \arq{tests/permissoes-de-iframe.spec.ts} falha se os dois primeiros
  divergirem: a divergência só aparece em produção, como 500 no
  \rota{PUT /experiencias/:expId}.\\
\cd{verbos} & \emph{objeto} & Documentação dos verbos de telemetria que a
  experiência emite. Não é lido por código; é contrato de leitura para quem
  analisa os dados.\\
\bottomrule
\end{tabularx}

\begin{regra}
\cd{series: []} significa \textbf{todas as séries}, não nenhuma: é assim que uma
experiência \emph{transversal} se declara (é o caso da app-70, que não pertence a
ano nenhum). Desde a sala de aula (\adrr{ADR-69}) este campo é \textbf{lido}: é ele
que decide se a missão entra na grade de uma turma, e vazio continuar querendo
dizer \enquote{todas} deixou de ser hipótese e virou o que sustenta a transversal
aparecer em toda sala. O contrário faz a experiência desaparecer da tela sozinha,
sem erro em lugar nenhum.

\textbf{Campo novo no \arq{experiencia.json} tem de ser acrescentado ao
\arq{publicar-manifesto.mjs} também} --- foi assim com \cd{permissoes}, com
\cd{unidade} e com \cd{produto}. Esquecer é INVISÍVEL em desenvolvimento (onde o
hub lê o arquivo do \emph{bundle}) e quebra em produção, onde ele lê a coluna. E
a cadeia tem CINCO pontos, não dois: o Zod da rota, a coluna, o
\arq{publicar-manifesto.mjs}, a \emph{interface} \cd{Manifesto} do \emph{seed}
(\arq{prisma/catalogo-e-liberacoes.ts}) e o \arq{\_template} --- este último
porque \cd{pnpm new:exp} copia o arquivo dele, e um \emph{template} incompleto
faz TODA experiência nova nascer errada.
\end{regra}

\subsection{Como o hub resolve a URL do \cd{iframe}}

\begin{figure}[H]
\centering
\fitwidth{%
\begin{tikzpicture}[x=1cm, y=1cm]

  % Vertical, em duas decisões: a 1ª decide de onde vêm viewport e
  % permissões; a 2ª, o endereço do bundle (apps/hub/lib/embed-experiencia.ts).
  \node[actstart] (ini) at (8,0) {};
  \node[actionsys, text width=74mm, anchor=north] (le) at (8,-0.7)
    {lê em paralelo, na api, com o \emph{token} recém-assinado:\\
     \texttt{GET /experiencias/:expId} (o manifesto) e\\
     \texttt{GET /recursos/efetivos} (a versão fixada na escola)};
  \node[decision] (d1) at (8,-3.55) {tem manifesto?};

  \node[action, text width=50mm] (vp1) at (3.3,-5.3)
    {\texttt{viewport} e \texttt{permissoes}\\vêm do MANIFESTO};
  \node[actionsys, text width=50mm] (vp2) at (12.7,-5.3)
    {lê \texttt{\{baseInterna\}/experiencia.json}\\
     (\texttt{EXP\_INTERNAL\_HOST} aplicado à base)};

  \node[decision] (d2) at (8,-7.75) {produção e com manifesto?};

  \node[action, text width=50mm] (com) at (3.3,-9.9)
    {\texttt{\{cdn\}/\{expId\}/\{versao\}/}\\
     \texttt{index.html?hubOrigin=\ldots}\\[2pt]
     \texttt{versao}: a fixada no escopo,\\senão a \texttt{versaoAtiva}};
  \node[actionsys, text width=50mm] (sem) at (12.7,-9.9)
    {\texttt{\{base\}/?hubOrigin=\ldots}\\[2pt]
     \texttt{base}: em dev, o \texttt{src} de \texttt{/dev-embed}\\
     ou a porta de \texttt{lib/experiencias-dev.ts};\\
     em produção, o CDN};

  \node[actionok, text width=58mm] (frame) at (8,-12.6)
    {\texttt{ExpFrame}: \emph{sandbox} + \texttt{allow},\\
     dentro de \texttt{DispositivoFrame} se houver \texttt{viewport}};
  \node[actend] (fim) at (8,-14.0) {};

  \coordinate (j1) at (8,-6.3);
  \coordinate (j2) at (8,-11.55);

  \draw[flow] (ini) -- (le);
  \draw[flow] (le) -- (d1);
  \draw[flow] (d1.west) -| (vp1.north) node[flowlbl, pos=0.25]{sim};
  \draw[flow] (d1.east) -| (vp2.north) node[flowlbl, pos=0.25]{não};
  \draw[line width=0.85pt, draw=unooInkSoft] (vp1.south) |- (j1) -| (vp2.south);
  \draw[flow] (j1) -- (d2.north);
  \draw[flow] (d2.west) -| (com.north) node[flowlbl, pos=0.25]{sim};
  \draw[flow] (d2.east) -| (sem.north) node[flowlbl, pos=0.25]{não};
  \draw[line width=0.85pt, draw=unooInkSoft] (com.south) |- (j2) -| (sem.south);
  \draw[flow] (j2) -- (frame.north);
  \draw[flow] (frame) -- (fim);

\end{tikzpicture}}
\caption[Resolução da URL do embed]{Resolução da URL do \emph{embed} e da casca
  de aparelho. Os dois caminhos convergem: mesmo sem manifesto, o hub obtém
  \cd{viewport} e \cd{permissoes} lendo o \arq{experiencia.json} servido ao lado
  do \arq{index.html} do \emph{bundle}. No caminho versionado, a versão é a
  fixada no escopo da escola (\adrr{ADR-87}), quando há, e senão a
  \cd{versaoAtiva} do catálogo.}
\label{fig:fluxo-embed}
\end{figure}

\begin{atencao}
A leitura do \arq{experiencia.json} do \emph{bundle} \textbf{não é só de
desenvolvimento}. Antes era, e uma \emph{stack} com a tabela \cd{experiencias}
vazia --- \emph{compose} local, ou ambiente novo antes do primeiro
\emph{publish} do CI --- servia a exp08ef1b \emph{esticada na janela} em vez de
dentro do aparelho de 375$\times$812. Pelo mesmo motivo, a app-70 abriria sem
\cd{allow-pointer-lock}: o jogo monta, a câmera não anda, e nada no console diz
por quê.
\end{atencao}

\begin{nota}
Essa leitura é \emph{server-side}, e \env{NEXT\_PUBLIC\_CDN\_EXP\_URL} é o
endereço do \textbf{navegador} --- de dentro de um contêiner,
\cd{localhost:5173} é o próprio contêiner. Daí as duas variáveis:
\env{EXP\_INTERNAL\_HOST} troca só o \emph{host} e \textbf{preserva a porta} (é
o que serve em desenvolvimento, onde cada experiência tem um \emph{dev server}
numa porta diferente), e \env{EXP\_INTERNAL\_URL} é um endereço fixo, que só
serve quando há \emph{uma} origem de \emph{bundles} --- a CDN, em produção.
\end{nota}

\subsection{Publicação e reversão}

\begin{tabularx}{\linewidth}{@{}L{34mm}X@{}}
\toprule
\headrow \thd{Operação} & \thd{Comando}\\
\midrule
Publicar (feito pelo CI) &
  \cd{gcloud storage rsync -r \{dir\}/dist gs://\{bucket\}/\{exp\}/\{ver\}/} +
  \cd{node apps/hub/scripts/publicar-manifesto.mjs \{dir\}}\\
\zebra Reverter &
  \cd{node apps/hub/scripts/publicar-manifesto.mjs <exp> --versao x.y.z}\\
\bottomrule
\end{tabularx}

O \emph{script} assina um \emph{token} de cinco minutos com a chave de assinatura
e chama \rota{PUT /experiencias/:expId} como \cd{admin-contrato}. A API valida
esse \emph{token} contra o JWKS que o \emph{hub} publica --- o que é a razão de o
\emph{workflow} publicar as experiências \emph{depois} do hub.

\section{Contrato de telemetria}

O formato do evento é inspirado em xAPI~\cite{xapi2023} e simplificado ao que a
plataforma efetivamente usa.

\begin{lstlisting}[language=jsonunoo, caption={Evento de telemetria (canal analítico)}, label={lst:evento}]
{
  "eventId":    "uuid",          // gerado pela SDK; chave de dedupe
  "ts":         "ISO-8601",      // relogio do cliente
  "estudanteId":    "...",           // SOBRESCRITO pelo servidor, a partir do token
  "escolaId":   "...",           // SOBRESCRITO pelo servidor
  "expId":      "exp08ef1b",
  "expVersion": "1.2.0",
  "sessaoId":   "uuid",
  "verbo":      "completed | attempted | interacted | scored | progressed | ...",
  "objeto":     "fase-3",
  "resultado":  { },             // livre, por experiencia
  "contexto":   { }              // livre, por experiencia
}
\end{lstlisting}

\begin{tabularx}{\linewidth}{@{}L{30mm}X@{}}
\toprule
\headrow \thd{Cláusula} & \thd{Regra}\\
\midrule
Quem preenche o quê & A SDK preenche tudo \textbf{exceto} \cd{verbo},
  \cd{objeto} e \cd{resultado} --- esses são a única coisa que a experiência
  informa.\\
\zebra Identidade & \cd{estudanteId} e \cd{escolaId} enviados pelo cliente são
  \textbf{descartados}: o servidor os deriva do \emph{token}. É anti-\emph{spoof}
  e, ao mesmo tempo, garantia de que o cliente não \emph{pode} colocar PII no
  lugar do id.\\
Lote & Até 50 eventos por requisição, \emph{flush} a cada
  \SI{10}{\second} ou quando o \emph{buffer} enche.\\
\zebra Fechamento de página & No \cd{pagehide}, \rota{POST /telemetry/beacon}
  com o \emph{token} \textbf{no corpo} --- \cd{sendBeacon} não envia cabeçalho.\\
Idempotência & \cd{ON CONFLICT (event\_id, ts) DO NOTHING}. O Pub/Sub entrega
  \emph{at least once}, então reentrega é rotina, não exceção.\\
\zebra Documentação dos verbos & Cada experiência declara seus \cd{verbos} no
  manifesto e no próprio README. Verbo não documentado é dado que ninguém
  consegue interpretar seis meses depois.\\
\bottomrule
\end{tabularx}

O detalhamento dos dois canais --- por que são dois, e o que acontece em cada um
--- é o Capítulo~\ref{cap:telemetria}.

\section{Contrato do \emph{provider} de IA}

\begin{lstlisting}[language=tsunoo, caption={\arq{apps/api/src/ai/provider.ts} --- a interface neutra de \emph{vendor}}, label={lst:aiprovider}]
export interface AiRequest {
  system: string;
  user: string;
  maxTokens?: number;
  /** Familia de modelo, NEUTRA de vendor. `padrao` e o modelo do dia a dia;
   *  `longo` e o de saida grande (planos anuais). Quem chama nao sabe qual e. */
  modelo?: "padrao" | "longo";
  temperatura?: number;
}

export interface AiResult {
  texto: string;
  tokensUsados: number;
  /** `STOP` (natural) ou outro motivo do vendor (`MAX_TOKENS`, `SAFETY`...).
   *  Aditivo: quem nao liga, ignora. Quem gera texto longo TEM de olhar. */
  finishReason?: string;
}

export interface AiFerramenta {
  nome: string;
  descricao: string;
  /** Reusa o AiJsonSchema neutro. `description` por PROPRIEDADE nao e
   *  comentario: e por ela que o modelo descobre o que pode mandar. */
  parametros: AiJsonSchema;
}

export interface AiConversaRequest {
  system: string;
  historico: Array<{ papel: "estudante" | "agente"; texto: string }>;
  ferramentas?: {
    declaracoes: AiFerramenta[];
    executar: (nome: string, args: Record<string, unknown>) => Promise<unknown>;
  };
  /** Motivos para RECUSAR a fala. Motivo derruba a tentativa e o modelo tenta
   *  de novo sabendo por que. As regras sao do dominio; o laco e do provider. */
  validar?: (texto: string) => string[];
  socorro?: string;              // o que pedir na ULTIMA tentativa
  maxVoltas?: number;            // teto de idas e voltas de ferramenta (6)
  tentativas?: number;           // default 3
  /** Opcionais (ADR-175): a chamada acima do teto e abortada e repetida (ate 3
   *  por volta); passado o prazo do turno, vira ErroDeGeracao. Sem eles, nada muda. */
  tetoPorChamadaMs?: number;
  prazoMs?: number;
}

export interface AiSpeechRequest {
  texto?: string;                                    // monologo
  turnos?: Array<{ falante: string; texto: string }>; // dialogo
  /** Timbre por falante, NEUTRO de vendor: `jovem-feminino`, `idoso-masculino`...
   *  Quem chama nunca escolhe a voz do catalogo do fornecedor. */
  falantes: AiFalante[];
  direcao?: string;              // atuacao: "tom assustado, voz tremula"
  idioma?: string;               // BCP-47, default pt-BR
}

export interface AiJsonRequest<T> {
  prompt: string;                // NAO `user`: o pedido estruturado tem nomes proprios
  system?: string;
  responseSchema: AiJsonSchema;  // NAO `schema`
  /** Valida e normaliza; lancar faz REGERAR com o motivo no prompt. */
  validar: (bruto: unknown) => T;
  temperatura?: number;
  tentativas?: number;
  /** Teto de saida. Ausente = o teto do modelo. Uma semana de aulas passa
   *  de 5k tokens, e JSON truncado nao chega como MAX_TOKENS legivel. */
  maxTokens?: number;
  modelo?: "padrao" | "longo";   // mesma familia neutra de AiRequest
  /** Quanto o vendor cobrou, UMA VEZ POR TENTATIVA -- inclusive as que
   *  `validar` rejeitou: foram pagas igual. Quem chama `temSaldo` TEM de
   *  passar isto; sem ele nao ha de onde tirar o numero (ADR-74).
   *  `divisao` (ADR-161): entrada e saida, quando o vendor informa; a quota
   *  ignora, o painel de custos usa. */
  aoUsarTokens?: (tokens: number, divisao?: DivisaoDeTokens) => void;
}

/** saida = total - entrada: inclui o raciocinio, cobrado na tarifa de saida. */
export interface DivisaoDeTokens { entrada: number; saida: number }

export interface AiImagemOpcoes {
  proporcao?: AiProporcao;
  /** O usageMetadata da geracao. So o MedidorDeProvider passa (ADR-161). */
  aoUsarTokens?: (tokens: number, divisao?: DivisaoDeTokens) => void;
}

export interface AiProvider {
  generate(req: AiRequest): Promise<AiResult>;
  /** Saida JSON estruturada, com retry e reforco do erro anterior no prompt. */
  generateJson<T>(req: AiJsonRequest<T>): Promise<T>;
  /** Bytes originais (jpeg/png). Falha => quem chama cai no fallback procedural.
   *  A proporcao e PEDIDA ao modelo (`opcoes.proporcao`), nunca recortada depois. */
  generateImage(prompt: string, opcoes?: AiImagemOpcoes): Promise<Buffer>;
  /** MP3 pronto para <audio src>. Falha => quem chama cai na transcricao. */
  generateSpeech(req: AiSpeechRequest): Promise<Buffer>;
  /** Conversa multi-turno com ferramentas (ADR-23). `tokensUsados` soma TODAS
   *  as voltas e TODAS as tentativas, inclusive as rejeitadas: e o que foi pago. */
  conversar(req: AiConversaRequest): Promise<AiResult>;
  /** Conversa SEM ferramentas, devolvida aos pedacos (ADR-132). */
  conversarEmFluxo(req: AiFluxoRequest): AsyncIterable<AiFluxoEvento>;
}

export interface AiFluxoRequest {
  system: string;
  historico: Array<{ papel: "pessoa" | "agente"; texto: string }>;
  maxTokens?: number;
  temperatura?: number;
  sinal?: AbortSignal;           // abortar corta a chamada ao vendor
  // sem `ferramentas` e sem `validar`: num stream o texto ja saiu
}

/** Pedacos de texto e, por ultimo, UM `fim` com o total -- tambem quando o
 *  modelo parou pelo teto. Erro de transporte e excecao, nao evento. */
export type AiFluxoEvento =
  | { tipo: "texto"; texto: string }
  | { tipo: "fim"; resultado: AiResult };

export const AI_PROVIDER = "AI_PROVIDER";   // token de injecao do Nest
\end{lstlisting}

\begin{regra}
\textbf{Nenhum módulo fora de \arq{src/ai} conhece o \emph{vendor}.} Trocar de
provedor é criar um arquivo novo (\arq{src/ai/<vendor>.provider.ts}) e registrá-lo
no módulo por \env{AI\_PROVIDER}. Implementações atuais:
\arq{vertex.provider.ts} (\emph{default}) e \arq{mock.provider.ts} (testes e
desenvolvimento sem GCP).

O corolário aparece em \cd{AiTimbre} e \cd{AiProporcao}: quem chama pede
\enquote{voz de senhora idosa} e \enquote{retrato 4:5}, e é o \emph{provider}
que traduz para o nome de voz e o \cd{imageConfig} do fornecedor do dia. Um id
de voz do catálogo do Google atravessando o domínio seria \emph{vendor} vazando
para dentro da missão.
\end{regra}

\begin{regra}
\textbf{Toda chamada de IA deixa uma linha de consumo, e quem chama não faz
nada para isso} (\adrr{ADR-161}). O \cd{MedidorDeProvider}, que já embrulhava o
\emph{provider} para as métricas, recebe o \cd{RegistroDeConsumo}
(\arq{ai/consumo-de-ia.ts}) e registra cada operação no fim dela, com o que foi
PAGO:
\begin{itemize}
  \item \cd{generate}, \cd{conversar} e \cd{conversarEmFluxo}: a entrada é o
    \cd{promptTokenCount}, e a saída é o total MENOS a entrada --- o que inclui
    o raciocínio, cobrado na tarifa de saída (Seção~\ref{sec:ia-raciocinio}).
  \item \cd{generateJson}: a soma das TENTATIVAS, inclusive as que \cd{validar}
    recusou, pelo \cd{aoUsarTokens(tokens, divisao)}. O segundo argumento é
    novo e opcional: quem só precisa do total (a quota) o ignora; sem divisão,
    o total conta como saída, a tarifa mais cara.
  \item \cd{generateImage}: pelo \cd{aoUsarTokens} novo de \cd{AiImagemOpcoes},
    que só o medidor passa, e uma imagem por geração que deu certo.
  \item \cd{generateSpeech}: \textbf{estimado}, porque a Cloud TTS não devolve
    contagem --- a duração sai do MP3 (32\,kbps), a saída de 25 \emph{tokens}
    por segundo e a entrada de $\approx$ 4 \emph{bytes} por \emph{token} do
    roteiro.
\end{itemize}
As duas portas fora do \cd{AiProvider} registram por conta própria: o agente do
professor no fim do turno (\cd{consumoDoTurno}, só eventos não parciais) e o
mascote uma linha por sessão, somando o \cd{usageMetadata} de cada turno. A
\textbf{atribuição} (escola, perfil e recurso) sai do contexto da requisição: o
\cd{contextoMiddleware} põe em \cd{ContextoDaRequisicao.recurso} o primeiro
segmento da rota, e a voz ao vivo, que chega por WebSocket fora desse contexto,
a recebe em \cd{PedidoDeVoz.atribuicao}. O registro \textbf{nunca} derruba a
chamada: uma falha de banco vira o \emph{log} \cd{ia-consumo-nao-gravado}.
\end{regra}

\begin{nota}
\cd{AiJsonSchema} é um subconjunto deliberadamente pequeno de JSON Schema
(\cd{object}, \cd{array}, \cd{string}, \cd{number}, \cd{boolean}, \cd{enum}), e
cada \emph{provider} o traduz para o formato do seu modelo. O campo
\cd{validar} do \cd{AiJsonRequest} recebe a saída bruta e \emph{lança} quando ela
não serve --- o que faz a tentativa ser regerada com o erro anterior reforçado no
\emph{prompt}. Validação como parte do contrato, não como verificação posterior.
\end{nota}

\begin{regra}
\textbf{\cd{generate} e \cd{generateJson} não compartilham os nomes dos campos},
e confundi-los não é erro de digitação: \cd{prompt}/\cd{responseSchema} contra
\cd{user}/\cd{schema}. Um \cd{as never} sobre o objeto literal cala o
compilador, e o pedido chega ao \emph{provider} sem \cd{responseSchema} ---
a montagem do corpo estoura \emph{antes} da requisição, as três tentativas
queimam em segundos e o professor vê \enquote{Falhou} com zero token gasto.
Foi o que aconteceu na app-92 (plano de aulas). \textbf{Nunca use \cd{as never}
para satisfazer este contrato}: é justamente ele que o typecheck existe para
conferir.

O segundo ensinamento é do laço de \emph{retry}: falha de \textbf{transporte}
não é falha de \textbf{conteúdo}. Reescrever o \emph{prompt} não devolve
permissão nem conserta projeto errado, então \cd{4xx} (menos \cd{429}) falha na
primeira tentativa e a mensagem preserva o \emph{status}; \cd{429} e \cd{5xx}
são passageiros e seguem no laço. Sem essa distinção o corpo de erro do
\emph{vendor} --- que não tem \cd{candidates} --- virava \enquote{resposta
vazia}, e a causa real morria atrás de uma mensagem genérica.
\end{regra}

\section{Contrato do agente de IA do professor}
\label{sec:contrato-agente}

O agente (\adrr{ADR-126}) é o único contrato da plataforma com \textbf{quatro
pontas}: o agente em Python (\arq{apps/agente-professor}), a api
(\arq{src/agente-professor} + \arq{src/ai/agent-engine.client.ts}), o hub (o chat
flutuante e a rota de \emph{stream}) e a infraestrutura (o repositório de
fundação). A fonte normativa é \arq{docs/regras/agente-professor.md}; esta seção
o resume para quem lê a arquitetura, e os \emph{payloads} estão no
Apêndice~\ref{ap:payloads}.

\begin{lstlisting}[language=tsunoo, caption={\arq{apps/api/src/ai/agente-professor.port.ts} --- a segunda porta de \arq{src/ai} (abreviado)}, label={lst:agenteport}]
export interface EventoAgente {        // um Event do ADK, NORMALIZADO
  autor: string;                       // "user" ou o nome do agente
  parcial: boolean;                    // partial=true do ADK em modo SSE
  textos: string[];                    // sem as partes de raciocinio
  chamadasDeFerramenta: Array<{ nome: string; args: Record<string, unknown> }>;
  respostasDeFerramenta: Array<{ nome: string; resposta: unknown }>;
  estadoDelta?: Record<string, unknown>;   // por aqui chega titulo_conversa
  tokens?: number;                     // usage_metadata; so vale nos NAO parciais
  erro?: { codigo: string; mensagem?: string };
}

export interface AgenteProfessorPort {
  readonly modo: "agent-engine" | "local" | "mock" | "desligado";
  criarSessao(usuarioId: string, estado: Record<string, unknown>,
              expiraEm?: Date): Promise<{ sessaoId: string }>;
  lerSessao(usuarioId: string, sessaoId: string): Promise<SessaoAgente>;
  apagarSessao(usuarioId: string, sessaoId: string): Promise<void>; // idempotente
  renovarSessao?(sessaoId: string, expiraEm: Date): Promise<void>;  // so agent-engine
  enviar(usuarioId: string, sessaoId: string, texto: string,
         signal: AbortSignal): AsyncIterable<EventoAgente>;
}
\end{lstlisting}

\begin{regra}
\textbf{Só \arq{agent-engine.client.ts} conhece o engine}, pela mesma regra do
\cd{AiProvider}. É porta própria, e não um método a mais no \cd{AiProvider},
porque modelo, ferramentas e instrução moram no agente: a api só conduz a
sessão. \textbf{Um cliente, dois alvos}: \cd{agent-engine} (\cd{:query} e
\cd{:streamQuery?alt=sse} do Agent Engine, com o ADC) e \cd{local} (o
\emph{shim} de \arq{apps/agente-professor}, que expõe O MESMO contrato em
\rota{/api/reasoning\_engine} e \rota{/api/stream\_reasoning\_engine}, sem
\emph{token}). \cd{mock} é o da suíte e2e; \cd{desligado} é o padrão.
\end{regra}

\subsection{O que a api pede ao engine}

\begin{tabularx}{\linewidth}{@{}L{40mm}X@{}}
\toprule
\headrow \thd{Chamada} & \thd{Uso}\\
\midrule
\cd{async\_create\_session} (\cd{:query}) & Cria a sessão com
  \cd{user\_id = users.id}, o estado \cd{\{nome\_professor\}} e
  \cd{expire\_time} distante pelos \cd{**kwargs}. Se o engine recusar o campo, a
  criação repete sem ele (TTL padrão de 365 dias).\\
\zebra \cd{async\_get\_session} (\cd{:query}) & Monta o histórico da tela (só o
  texto do professor e o texto FINAL do agente) e lê o título quando o
  \emph{stream} terminou sem ele. \enquote{Não existe} vira 404 e apaga a linha.\\
\cd{async\_delete\_session} (\cd{:query}) & Apagar. Idempotente: sessão que já
  não existe não é erro.\\
\zebra \cd{async\_stream\_query} (\cd{:streamQuery?alt=sse}) & O turno, com
  \cd{run\_config.streaming\_mode = "sse"}. A resposta vem em linhas \cd{data:}
  \emph{ou} NDJSON --- o cliente decide por linha, e linha que não é JSON é
  contada e vai a log, sem derrubar o turno.\\
\cd{PATCH v1beta1/\ldots/sessions/\{id\}?\B updateMask=expireTime} & A renovação
  depois de cada turno entregue. É a API NATIVA: o \cd{:query} do ADK não tem
  método de atualizar sessão, e a SA da api precisa de \cd{sessions.update}.
  Falha é só log (\cd{agente-professor.renovacao-falhou}).\\
\bottomrule
\end{tabularx}

\subsection{As chaves de estado e as ferramentas}

\begin{tabularx}{\linewidth}{@{}L{30mm}L{30mm}X@{}}
\toprule
\headrow \thd{Chave} & \thd{Quem escreve} & \thd{Quem lê}\\
\midrule
\cd{nome\_professor} & a api, na criação (primeiro nome de
  \cd{users.display\_name}) & a instrução do agente\\
\zebra \cd{titulo\_conversa} & \textbf{o agente}, com o próprio LLM, depois do
  primeiro turno & a api (no \cd{actions.state\_delta} do \emph{stream}, ou numa
  leitura da sessão) e o hub\\
\bottomrule
\end{tabularx}

As duas ferramentas são de \emph{data store} (\adrr{ADR-127}).
\cd{buscar\_habilidades} consulta o estruturado com filtro exato e paginação;
\cd{buscar\_acervo} devolve trechos dos PDFs, e as chaves de cada trecho ---
\cd{titulo}, \cd{pagina\_inicio}, \cd{pagina\_fim}, \cd{trecho}, \cd{fonte\_id}
--- \textbf{são lidas pela api} para o evento \cd{fontes}. Renomear uma delas sem
mudar a api junto faz as fontes sumirem da tela, sem erro.

\subsection{O que a api emite ao navegador}

\rota{POST /agente-professor/conversas/:id/mensagens} \cd{\{ texto \}} (1 a 4.000
caracteres) responde \cd{text/event-stream} (\adrr{ADR-128}):

\begin{tabularx}{\linewidth}{@{}L{18mm}X@{}}
\toprule
\headrow \thd{\cd{event:}} & \thd{\cd{data:}}\\
\midrule
\cd{delta} & \cd{\{ texto \}} --- dos eventos parciais, e do final agregado SÓ no
  que os parciais não cobriram (a recusa de escopo chega num evento final sem
  parcial antes, com \cd{custom\_metadata.fora\_do\_escopo = true}).\\
\zebra \cd{fontes} & \cd{\{ fontes: [\{ titulo, paginaInicio, paginaFim \}] \}}, dos
  retornos de \cd{buscar\_acervo}, sem repetição.\\
\cd{titulo} & \cd{\{ titulo \}}, só quando muda.\\
\zebra \cd{aviso} & \cd{\{ codigo: "habilidade-inexistente", mensagem \}}, um por
  código citado que \cd{habilidadePorCodigo} não conhece.\\
\cd{fim} & \cd{\{\}}.\\
\zebra \cd{erro} & \cd{\{ codigo, mensagem \}}, com \cd{codigo} em \cd{QUOTA},
  \cd{RATE\_LIMIT}, \cd{AGENTE\_INDISPONIVEL} ou \cd{TEMPO\_ESGOTADO}.\\
\bottomrule
\end{tabularx}

Mais \cd{: ping} a cada 15\,s. O \emph{stream} termina em \cd{fim} \textbf{ou} em
\cd{erro}; fontes, título, avisos e \cd{fim} saem depois do último evento do
agente, nessa ordem. Tudo o que recusa responde JSON \textbf{antes} de abrir o
\emph{stream}.

\begin{atencao}
\textbf{Três armadilhas medidas, que nenhum tipo pega.}
\begin{enumerate}
  \item \textbf{O ADK em SSE emite os parciais E um evento final agregado que
    REPETE o texto inteiro.} Repassar os dois duplica a resposta na tela; o
    \cd{DeduplicadorDeTexto} deixa passar do final só o que faltou, por segmento
    (uma chamada ao modelo).
  \item \textbf{\cd{:streamQuery?alt=sse} responde ora \cd{data:}, ora NDJSON}
    (\emph{issue} aberta no \cd{python-aiplatform}). O \emph{parser} aceita os dois,
    por linha.
  \item \textbf{\emph{Tokens} só dos eventos NÃO parciais}, somados: um turno com
    ferramenta faz mais de uma chamada ao modelo, e os parciais repetem o
    \cd{usage\_metadata}. A guarda de escopo e o título somam os deles ao
    primeiro evento não parcial seguinte.
\end{enumerate}
\end{atencao}

\textbf{Política de evolução.} O \cd{ContratoInterceptor} não confere o SSE: quem
confere cada \cd{data:} é o teste, contra \cd{DADOS\_POR\_EVENTO} de
\arq{contrato-de-rotas.ts}, e \textbf{evento novo entra lá antes de sair da api}.
O ADK está fixado no \arq{uv.lock} (\cd{\textasciitilde=2.10.0}); subir de versão
é refazer as medições de \arq{tests/test\_e2e.py} do agente, porque é delas que
saem as três armadilhas acima. Chave de estado, nome de ferramenta e chave de
trecho mudam nas quatro pontas no mesmo PR.

\section{Contrato com o \emph{vendor} do avatar da Ajuda}
\label{sec:contrato-ajuda}

A Ajuda com avatar (\adrr{ADR-132}, Seção~\ref{sec:ia-ajuda}) é o único contrato
em que a api \textbf{é chamada pelo servidor de um \emph{vendor}} --- e não por
uma pessoa logada nem por um serviço do Google. São três trechos: a api pede a
sala ao LiveAvatar; o navegador entra na sala com o SDK do \emph{vendor}
(WebRTC); e, a cada fala da pessoa, o \emph{vendor} chama o LLM da api, num
formato OpenAI-compatível que \textbf{a documentação dele não descreve} --- tudo
o que segue sobre esse trecho foi MEDIDO no \emph{spike} de 2026-10-02, no
\emph{sandbox}. A doutrina inteira está em \arq{apps/api/src/ajuda/AGENTS.md}; os
\emph{payloads}, na Seção~\ref{sec:payload-ajuda}.

\begin{lstlisting}[language=tsunoo, caption={\arq{apps/api/src/ai/avatar-ao-vivo.port.ts} --- a terceira porta de \arq{src/ai} (abreviado)}, label={lst:avatarport}]
export const AVATAR_AO_VIVO = "AVATAR_AO_VIVO";

export interface PedidoDeSessao {
  duracaoMaximaS: number;              // o vendor derruba a sala ao atingi-lo
  variaveis: Record<string, string>;   // dynamic_variables: por aqui vai o token
}

export interface SessaoDeAvatar {
  sessaoId: string;                    // o id da sala no vendor (para encerrar)
  tokenDaSessao: string;               // o que o NAVEGADOR usa para entrar
}

export interface AvatarAoVivoPort {
  readonly ligado: boolean;            // false = AvatarDesligado: a Ajuda da 503
  criarSessao(pedido: PedidoDeSessao): Promise<SessaoDeAvatar>;
  encerrarSessao(sessaoId: string): Promise<void>;   // sala ja fechada nao e erro
}
\end{lstlisting}

\begin{regra}
\textbf{Só \arq{liveavatar.client.ts} conhece o LiveAvatar}, pela regra do
\cd{AiProvider}; \arq{src/ajuda} fala com a porta. É porta própria, e não método
do \cd{AiProvider}, porque não gera texto: reserva uma sala cobrada por minuto,
com ciclo de vida. O texto da Ajuda passa pelo \cd{AiProvider} como qualquer
outro, pelo \cd{conversarEmFluxo} (Listagem~\ref{lst:aiprovider}).
\end{regra}

\subsection{O que a api pede ao \emph{vendor}}

\begin{tabularx}{\linewidth}{@{}L{40mm}X@{}}
\toprule
\headrow \thd{Chamada} & \thd{Uso}\\
\midrule
\rota{POST /v1/sessions/token} & Abre a sala: \cd{mode: "FULL"},
  \cd{interactivity\_type: "CONVERSATIONAL"}, o avatar, o teto
  (\cd{max\_session\_duration}), \cd{llm\_configuration\_id} (o NOSSO LLM),
  \cd{avatar\_persona} com \cd{context\_id} e \cd{language: "pt"}, e as
  \cd{dynamic\_variables} com o \emph{token} da sessão. Devolve
  \cd{session\_id} e \cd{session\_token}; só o segundo vai ao hub.\\
\zebra \rota{POST /v1/sessions/stop} & Fecha a sala (\cd{reason: "USER\_CLOSED"}).
  4xx é sala que já fechou, e é ignorado.\\
\bottomrule
\end{tabularx}

Cabeçalho \cd{x-api-key} com \env{LIVEAVATAR\_API\_KEY}, 10\,s de \emph{timeout}
--- abrir é o passo em que a pessoa espera olhando um vídeo preto. O contexto e a
configuração de LLM que a sessão cita são criados uma vez por contrato e
ambiente por \arq{scripts/liveavatar-configurar.mjs} (idempotente pelo nome
\cd{unooria-ajuda-<ambiente>}), \textbf{fora do Terraform}.

\subsection{O que o \emph{vendor} pede à api}

\begin{tabularx}{\linewidth}{@{}L{40mm}X@{}}
\toprule
\headrow \thd{Chamada} & \thd{O que chega, e o que a api responde}\\
\midrule
\rota{GET /ajuda/llm/models} & Na abertura da sala, com 5\,s de paciência.
  Resposta mínima: um modelo, \cd{unooria-ajuda}. O modelo de verdade
  (\env{VERTEX\_MODEL}) não é anunciado.\\
\zebra \rota{POST /ajuda/llm/chat/\B completions} & A cada fala. É o \emph{request}
  do cliente OpenAI em Python (\emph{LiveKit Agents}):
  \cd{Authorization: Bearer <segredo>} e
  \cd{\{messages, model, stream: true, stream\_options: \{include\_usage: true\},
  temperature: 0.5\}}. \cd{messages[0]} é \cd{system} = o \emph{prompt} do
  contexto com as \cd{dynamic\_variables} JÁ substituídas; depois vem o
  histórico inteiro, a começar pelo \cd{opening\_text}. Resposta:
  \cd{chat.completion.chunk} em SSE, ou um \cd{chat.completion} inteiro quando
  \cd{stream} é falso.\\
\bottomrule
\end{tabularx}

\begin{regra}
\textbf{Duas trancas, e as duas são necessárias.} A rota fica fora do
\cd{OidcGuard} --- quem chama é o \emph{vendor} ---, e no lugar dele exige:
\begin{enumerate}
  \item \cd{Authorization: Bearer} \env{AJUDA\_LLM\_SEGREDO}, comparado em tempo
    constante: prova que é o \emph{vendor}. No \emph{vendor}, o mesmo valor mora
    como segredo do tipo \cd{OPENAI\_API\_KEY} da nossa configuração de LLM;
  \item o \emph{token} da sessão \textbf{dentro da mensagem \cd{system}}, no
    marcador \cd{[sessao:<token>]}: prova QUEM fala, e numa sessão que esta api
    abriu e não fechou (\cd{ajuda\_sessoes.fim} nulo). É um HS256 curto
    (\env{AJUDA\_SESSAO\_CHAVE}; \cd{\{sub, escola, perfil, sid\}}; audiência
    \cd{ajuda-llm}; vale o teto da sessão mais 1 minuto) que a api assina ao
    abrir e entrega em \cd{dynamic\_variables.sessao}. HS256, e não o RS256 do
    hub, porque quem assina e quem confere é a mesma api.
\end{enumerate}
Só o segredo deixaria qualquer um que o tivesse conversar em nome de ninguém; só
o \emph{token}, qualquer um com um \emph{token} vazado usar o modelo direto.
\emph{Token} escrito na fala da pessoa não vale: só o \cd{system} é lido. O
\emph{token} não leva nome nem e-mail, porque o \emph{vendor} vê o \emph{prompt}
do contexto (regra 6).
\end{regra}

\begin{atencao}
\textbf{Quatro comportamentos medidos que a documentação do \emph{vendor} não
diz.}
\begin{enumerate}
  \item \textbf{Erro do endpoint = avatar MUDO}: sem nova tentativa e sem evento
    para o navegador. Por isso, aceita a chamada, \textbf{toda recusa de negócio
    e toda falha do modelo viram FALA, em 200}; só a autenticação responde 401
    (e a rota responde 503 sem as duas \emph{envs}).
  \item \textbf{\emph{Stream} é obrigatório na prática, e o primeiro byte tem
    10\,s} (\cd{x-stainless-read-timeout}). Daí o \cd{conversarEmFluxo}, o
    pedaço vazio com \cd{role: "assistant"} logo ao abrir o \emph{stream} e a
    fala de espera aos 7\,s.
  \item \textbf{Sem \cd{context\_id} na sessão, o avatar fica mudo}, sem erro.
  \item \textbf{O segredo do LLM é \cd{secret\_type: OPENAI\_API\_KEY}}: a
    \emph{skill} da HeyGen ensina \cd{LLM\_API\_KEY}, que a api deles recusa.
\end{enumerate}
\end{atencao}

\textbf{Política de evolução.} O corpo é lido com \cd{.passthrough()}: só
\cd{messages} decide, e campo novo do \emph{vendor} é tolerado --- recusar com 400
deixaria o avatar mudo sem erro na tela. O \cd{ContratoInterceptor} não confere a
resposta (é \cd{text/event-stream}, escrita crua). Mudou o formato do
\emph{request}, a substituição das \cd{dynamic\_variables} ou o \emph{timeout} do
lado do \emph{vendor}? O desenho inteiro sai dessas medições: refaça o
\emph{spike} no \emph{sandbox} antes de mexer na api. Trocar a URL pública da api
ou o segredo pede recriar a configuração de LLM no \emph{vendor} (ele não tem
edição documentada; apagar o segredo ANTES da configuração a quebra sem aviso).

\section{Contrato da conexão de voz do mascote da Ajuda}
\label{sec:contrato-ajuda-mascote}

O mascote (\adrr{ADR-158}) é o segundo atendente da Ajuda, ao lado do avatar.
Aqui o \emph{vendor} é o Gemini Live da Vertex, uma conversa de VOZ bidirecional,
e o rosto é nosso: o mascote 3D do hub, com a boca tirada do próprio áudio. São
dois trechos, e a api fica no meio dos dois: \textbf{o navegador fala com a api}
por \rota{WS /ajuda/mascote/ao-vivo}, num protocolo NEUTRO e fechado; \textbf{só
a api fala com o Gemini}, com o ADC, pela porta \cd{VOZ\_AO\_VIVO}. Nada do
formato do \emph{vendor} atravessa --- nem para lá, nem para cá ---, e a
instrução do sistema (a mesma base da Ajuda, com a apresentação do mascote) é
montada na api e nunca chega ao navegador (regra 2).

\begin{lstlisting}[language=tsunoo, caption={\arq{apps/api/src/ai/voz-ao-vivo.port.ts} --- a quarta porta de \arq{src/ai} (abreviado)}, label={lst:vozport}]
export const VOZ_AO_VIVO = "VOZ_AO_VIVO";

export interface PedidoDeVoz {
  modelo: string;      // do painel (ajuda_configuracao), ex.: gemini-3.8-live
  voz: string;         // uma das vozes pre-definidas do Gemini, ex.: Leda
  instrucao: string;   // o prompt do atendente, com a base: nunca vai ao hub
  /** A quem o painel de custos atribui a sessao (ADR-161): o WebSocket chega
   *  por `upgrade`, fora do middleware que abre o contexto da requisicao. */
  atribuicao?: { escolaId?: string | null; perfil?: string | null; recurso?: string };
}

export type EntradaDaVoz =             // o que o navegador pode mandar
  | { tipo: "texto"; texto: string }
  | { tipo: "audio"; base64: string }  // 100 ms de PCM 16 bits, 16 kHz
  | { tipo: "fim-do-audio" };

export type EventoDaVoz =              // o que volta ao navegador
  | { tipo: "audio"; base64: string; taxa: number }   // PCM 16 bits, 24 kHz
  | { tipo: "transcricao-saida"; texto: string }
  | { tipo: "transcricao-entrada"; texto: string }
  | { tipo: "interrompido" }           // falou por cima: descartar a fila
  | { tipo: "fim-do-turno" };

export interface VozAoVivoPort {
  readonly ligada: boolean;            // false sem projeto da Vertex (mock)
  abrir(pedido: PedidoDeVoz): Promise<ConversaDeVoz>;
}
\end{lstlisting}

\begin{regra}
\textbf{Só \arq{gemini-live.client.ts} conhece o \cd{BidiGenerateContent}}, pela
regra do \cd{AiProvider}: \arq{src/ajuda} fala com a porta, e trocar de
\emph{vendor} é escrever outro cliente. É porta própria pela razão da do avatar:
não gera texto, abre uma conversa com ciclo de vida. O cliente abre
\cd{wss://<regiao>-aiplatform.googleapis.com/ws/\B google.cloud.aiplatform.\B
v1beta1.\B LlmBidiService/\B BidiGenerateContent} com o ADC, na região
\env{AJUDA\_MASCOTE\_REGIAO} (o Live não está no \cd{global}), e o modelo vai com
o caminho inteiro \cd{projects/<p>/locations/<r>/publishers/google/models/<m>}.
\end{regra}

\subsection{O protocolo entre o navegador e a api}

\begin{tabularx}{\linewidth}{@{}L{34mm}L{18mm}X@{}}
\toprule
\headrow \thd{Mensagem} & \thd{Sentido} & \thd{Regra}\\
\midrule
\cd{\{tipo: "entrar", bilhete\}} & \textbf{1ª}, ida & SEMPRE a primeira. O
  bilhete é o que \rota{POST /ajuda/sessoes} devolveu: HS256 com
  \env{AJUDA\_SESSAO\_CHAVE}, audiência \cd{ajuda-mascote}, 2 minutos, claims
  \cd{\{sub, escola, perfil, sid\}}. Vai na mensagem, e não na URL, que acaba em
  \emph{log} de balanceador. \textbf{Uma entrada por bilhete}: a api marca a
  linha de \cd{ajuda\_sessoes} de forma atômica. Sem bilhete em 10\,s, fecha.\\
\zebra \cd{\{tipo: "texto", texto\}} & ida & O que a pessoa escreveu, de 1 a 500
  caracteres.\\
\cd{\{tipo: "audio", base64\}} & ida & 100\,ms de microfone, PCM 16 bits
  \emph{little-endian}, 16\,kHz.\\
\zebra \cd{\{tipo: "fim-do-audio"\}} & ida & O microfone desligou: o modelo fecha o
  que ouviu e responde.\\
\cd{\{tipo: "pronto", prazo\}} & volta & O Gemini confirmou a abertura; \cd{prazo}
  é o fim da sessão (início + teto), em ISO.\\
\zebra \cd{\{tipo: "audio", base64, taxa\}} & volta & A fala do mascote, PCM 16 bits
  a 24\,kHz. É dela que o hub tira a boca.\\
\cd{transcricao-saida}, \cd{transcricao-entrada} & volta & O texto do que o
  mascote disse e do que a pessoa disse, para a conversa escrita da janela.\\
\zebra \cd{interrompido}, \cd{fim-do-turno} & volta & A pessoa falou por cima (o
  hub descarta o áudio na fila); o turno do mascote acabou.\\
\bottomrule
\end{tabularx}

Mensagem que não é JSON, fora do \cd{MensagemDoNavegadorSchema} (Zod,
\cd{.strict()}) ou acima de 30 por segundo fecha a conexão. Os fechamentos são o
canal de erro, porque WebSocket aberto não tem \emph{status}:

\begin{tabularx}{\linewidth}{@{}L{16mm}X@{}}
\toprule
\headrow \thd{Código} & \thd{Quando}\\
\midrule
4401 & O bilhete não chegou em 10\,s, não confere, venceu ou já foi usado.\\
\zebra 4403 & O produto \cd{ajuda} está desligado para a escola --- na entrada ou
  numa conferência.\\
4408 & O teto da sessão chegou, ou a sessão foi encerrada (outra aba,
  \rota{POST /ajuda/sessoes/:id/encerrar}, outra sessão aberta).\\
\zebra 4429 & Mensagem demais, ou fora do protocolo.\\
4502 & O Gemini não abriu a sessão, ou caiu no meio.\\
\zebra 4503 & A stack não tem a Vertex nem a \env{AJUDA\_SESSAO\_CHAVE}.\\
\bottomrule
\end{tabularx}

\begin{regra}
\textbf{O custo tem teto, e quem o aplica é a api.} A conexão cai no prazo da
sessão (o mesmo cálculo de minutos do avatar), e a cada 30\,s a api confere se a
linha de \cd{ajuda\_sessoes} continua aberta e se o produto continua ligado:
encerrar em outra aba, ou o \cd{admin} desligar a Ajuda, derruba a conversa no
ar. Banco fora do ar por um instante não derruba: o prazo ainda vale. A queda da
conexão, de qualquer lado, grava o \cd{fim} da linha.
\end{regra}

\begin{atencao}
\textbf{O \cd{Origin} é conferido à mão}, antes do \emph{upgrade}, contra a mesma
lista do CORS (\env{CORS\_ORIGINS}, \arq{src/config/origens.ts}): o CORS não
alcança WebSocket. Fora da lista, 403. Não é a tranca principal --- sem bilhete
ninguém passa da primeira mensagem --- mas fecha a porta a qualquer site aberto no
navegador. \textbf{O que foi dito ou escrito não vai ao \emph{log}} (regra 6): a
linha da sessão é \cd{ajuda.mascote}, com \cd{sid}, \cd{perfil}, \cd{duracaoMs},
\cd{turnos} e \cd{motivo}.
\end{atencao}

\textbf{Três comportamentos do \emph{vendor} medidos no PoC} (que a api herdou):
\cd{responseModalities} e \cd{speechConfig} vão DENTRO de
\cd{generationConfig}, e a \cd{systemInstruction} tem de ser um \cd{Content}
(texto puro é recusado); o servidor valida o JSON da abertura antes da
credencial, então campo errado volta 1007 \enquote{Unknown name} até com
credencial ruim; e chega \cd{sessionResumptionUpdate} sem pedir, que é
ignorado. \textbf{Política de evolução}: mensagem nova do protocolo entra no
\cd{MensagemDoNavegadorSchema} e em \cd{EventoDaVoz} antes de sair do hub ou da
api, e \arq{tests/ajuda-mascote.spec.ts} prende os dois lados.

\section{Contrato entre duas experiências}

Quando uma experiência precisa do resultado de outra, a comunicação nunca é
direta --- módulos de domínio não se conhecem (Seção~\ref{sec:c4l3}). O contrato
é uma tabela nomeada com os dois números e um \emph{endpoint} de fronteira.

\begin{figure}[H]
\centering
\fitwidth{%
\begin{tikzpicture}[x=1cm, y=1cm]

  % Linha de cima: o jogo entrega o dado bruto; o professor dispara o colapso,
  % que grava a tabela de contrato. Linha de baixo: o Instanooria a lê.
  \node[c4person, text width=22mm, minimum height=10mm] (prof) at (7.0,2.7)
    {\textbf{Professor}};
  \node[bundle, text width=28mm] (v) at (1.7,0) {\textbf{app-70}\desc{Vila do Lume}};
  \node[c4component, text width=34mm] (mv) at (7.0,0)
    {\texttt{vila/quadrantes.ts}\desc{módulo \texttt{vila} da api:\\colapso 10 \seta{} 4}};
  \node[c4component, text width=38mm, fill=unooOrangeWash, draw=unooOrange] (t) at (12.6,0)
    {\texttt{exp08ef1b\_app\_70\_}\\\texttt{perfis\_importados}\desc{a tabela de contrato}};
  \node[c4component, text width=38mm] (mi) at (12.6,-3.2)
    {\texttt{perfis/largada.ts}\desc{módulo \texttt{perfis} da api:\\decide o lado do \emph{feed}}};
  \node[bundle, text width=28mm] (i) at (7.0,-3.2) {\textbf{exp08ef1b}\desc{Instanooria}};

  \draw[c4rel] (v) -- (mv) node[cxlabel, midway, above=1pt]{\emph{payload}\\bruto};
  \draw[c4rel] (prof) -- (mv) node[cxlabel, midway, right=2pt, align=left]
    {\texttt{POST /atividades/:id/}\\\texttt{perfis/da-vila}};
  \draw[c4rel] (mv) -- (t) node[cxlabel, midway, above=1pt]{grava\\(\emph{upsert})};
  \draw[c4rel] (t) -- (mi) node[cxlabel, midway, right=2pt]{lê};
  \draw[c4rel] (mi) -- (i) node[cxlabel, midway, above=1pt]{lado inicial};

  \node[actnote, text width=150mm, anchor=north west] at (0.15,-4.6)
    {\textbf{Duas travas que custam estudantes, de propósito:} valor não observado
     (\texttt{indicative} ou \texttt{none}) \textbf{fica fora da média} --- nunca
     entra como zero; e quadrante \emph{sem nenhum membro observado} invalida o
     perfil inteiro. Sem perfil, o Instanooria cai no sorteio balanceado, que é
     honesto sobre não saber. O preço é real: parte da turma não terá perfil
     importado.};

\end{tikzpicture}}
\begin{legendadiag}
\amostra{bundle} experiência (\emph{bundle}) \quad
\amostra{c4component} módulo da api \quad
\amostra{c4component, fill=unooOrangeWash, draw=unooOrange} a tabela de contrato \quad
\amostra{fill=unooNavy, draw=unooNavy, rounded corners=1pt} pessoa
\end{legendadiag}
\caption[Contrato entre duas experiências]{O único acoplamento entre
  experiências, e como ele é explícito: uma tabela cujo nome carrega os dois
  números, escrita por um \emph{endpoint} de fronteira acionado pelo professor.}
\label{fig:contrato-exps}
\end{figure}

\begin{regra}
O disparo é do \textbf{professor}, e \cd{lado\_forcado} é \textbf{preservado} no
\emph{upsert}. Importar sozinho --- ou levar o objeto inteiro no \cd{update} ---
desfaria em silêncio uma redistribuição de turma que o professor fez antes da
aula. \env{PERFIL\_IMPORT\_TOKEN} e \rota{POST /perfis/importar} continuam vivos
apenas como compatibilidade com a ponte externa original.
\end{regra}
